\documentclass{article}
\IfFileExists{neurips_2026.sty}{
  \usepackage[dblblindworkshop]{neurips_2026}
}{
  \usepackage[margin=1in]{geometry}
  \usepackage[round,authoryear]{natbib}
}

\usepackage{amsmath,amssymb,amsthm,booktabs,microtype,multirow}
\usepackage{graphicx}
\usepackage[hidelinks]{hyperref}
\usepackage{xcolor}
\usepackage{xspace}

\newcommand{\method}{\textsc{RAVS}\xspace}
\newcommand{\risk}{\mathcal{R}}
\newcommand{\data}{\mathcal{D}}
\newcommand{\family}{\Lambda}
\newcommand{\indicator}{\mathbf{1}}
\newcommand{\nfe}{\operatorname{NFE}}
\newcommand{\pending}[1]{\fcolorbox{orange!70!black}{orange!8}{\parbox{0.92\linewidth}{#1}}}
\newtheorem{proposition}{Proposition}

\title{Lossless Risk-Adaptive Macro-Speculation for Diffusion Language Models}
\author{Anonymous Authors}
\date{}

\begin{document}
\maketitle

\begin{abstract}
Aggressive early stopping reduced neural function evaluations (NFEs) in our diffusion language
model decoders, but changed answers catastrophically on Dream; conservative stopping preserved
answers but saved little. We replace stopping with \emph{Risk-Adaptive Verified Speculation}
(\method). The model cheaply drafts several future AdaBlock refinement states, evaluates the root
and all drafts as one batch, and traverses a draft edge only when the verified native AdaBlock
transition equals its complete child state. A mismatch returns the verified successor, never the
draft. Consequently, under deterministic batch-invariant inference, \method reproduces the entire
AdaBlock state trajectory and output exactly, regardless of router error. A learned local-risk
score controls only speculative depth---hence wasted batch work---and has no authority over token
acceptance. Unlike prior lossless self-speculation for one-token stepwise decoding, our verifier
speculates over AdaBlock's threshold-parallel, delimiter-adaptive macro-transition. We freeze three
depth routers and evaluate LLaDA-8B-Instruct and Dream-7B-Instruct on reasoning and code with exact
sequence audits, physical NFE, verifier-node work, and wall-clock latency. A positive claim requires
zero sequence disagreements and a paired-bootstrap NFE-reduction lower bound above 5\% for each
model. Numerical claims are inserted only from complete non-mock artifacts that pass these gates.
\end{abstract}

\section{Introduction}

Masked diffusion language models (dLMs) generate many tokens in parallel by repeatedly
denoising a partially masked sequence \citep{nie2025llada,ye2025dream}. Their wall-clock
advantage depends on when tokens or blocks are committed. Fixed schedules spend refinement
steps uniformly; recent adaptive methods vary the block length, token transfer, or stopping
time using confidence and stability signals
\citep{lu2026adablock,wu2025fastdllm,israel2025apd,mohamed2025sched}.

The central difficulty is that confidence is useful but not safe. In our development runs, a local
risk estimator achieved AUROC above $0.98$ on both model families, yet using its score as permission
to terminate refinement produced a 16--32 point accuracy loss on Dream. Making the gate strict
removed the loss but also reduced savings to roughly 1\%. This is a structural conflict: a fallible
predictor cannot simultaneously provide a strong accuracy guarantee and aggressively change the
trajectory.

Our answer changes what prediction is allowed to control. \method uses the risk score only to
allocate parallel verification capacity. It drafts a chain of plausible future AdaBlock states,
batches them through the same model, and accepts only states proven to lie on the reference
trajectory. Router error can waste parallel work, but cannot change a token. The problem therefore
becomes maximizing exact draft acceptance subject to a hardware budget, rather than predicting
when it is safe to abandon the reference decoder.

Our contributions are: (i) a lossless speculative verifier for AdaBlock's threshold-parallel,
delimiter-adaptive state transition; (ii) a formal trajectory-equivalence result that makes output
safety independent of risk-model calibration; (iii) a risk-adaptive capacity rule that trades
verification width for acceptance without granting commitment authority; and (iv) an auditable
two-model evaluation that reports physical calls, batched node work, latency, exact output parity,
and frozen compute gates.

\section{Related work}

\paragraph{Efficient dLM decoding.}
LLaDA and Dream establish competitive autoregressive-style reasoning with masked diffusion
objectives \citep{nie2025llada,ye2025dream}. Fast-dLLM uses confidence-aware parallel decoding
and cache reuse \citep{wu2025fastdllm}; AdaBlock adapts block size online
\citep{lu2026adablock}; APD learns token-wise denoising paths \citep{israel2025apd}; and SchED
uses similarity-based early stopping \citep{mohamed2025sched}. FreeDave, OSDT, and WINO
further explore lossless verification, one-shot thresholding, and revocable inference
optimization \citep{liu2026freedave,zhou2026osdt,wang2026wino}. DiCo and SOAR study
distillation and schedule optimization, respectively \citep{luo2026dico,cao2026soar}.
Very recent work predicts semantic block boundaries (SemBlock), combines confidence with
token and distribution persistence (LESS), or coordinates simultaneous commits through a
mean-field objective \citep{song2026semblock,mohamed2026less,zoabi2026meanfield}.
Stability-Weighted Decoding gives a temporal-distribution view of unsafe commits, STDec combines
spatial and temporal consistency, and Just-on-Time studies independent token convergence
\citep{wu2026swd,chen2026stdec,kohut2026jit}. Prophet commits early from answer-region confidence,
while the concurrent LATCH framework verifies sustained stability over a parsed candidate span
\citep{li2025prophet,lee2026latch}. We therefore use temporal JS only as an input feature and do
not claim stability itself as our contribution. Concurrent theory analyzes efficiency
conditions for confidence decoding \citep{cai2026confidence}; our guarantee instead follows from
exact transition verification. A contemporaneous
re-evaluation also finds that confidence-based remasking gains can depend strongly on the
decoding setting \citep{frkovic2026reevaluating}, reinforcing the need for paired,
protocol-specific evidence.

\paragraph{Self-speculative decoding.}
SSD first showed lossless dLM self-speculation by constructing a linear hierarchical verification
tree for one-token stepwise decoding \citep{gao2025ssd}. S2D2 instead routes between diffusion
drafting and an autoregressive mode of the same block-diffusion model \citep{han2026s2d2}.
Our work builds directly on SSD's lossless verification principle, but changes both the reference
operator and the capacity rule. A verified edge in \method is a complete AdaBlock macro-step that
may reveal several threshold-qualified tokens, while block boundaries remain delimiter-adaptive;
a learned trajectory-risk score chooses depth but never acceptance. Thus the relevant comparison
is fixed-depth SSD versus risk-adaptive speculation over an already adaptive parallel decoder.

\section{Risk-adaptive verified macro-speculation}

\subsection{The native AdaBlock transition}

Fix one active block and let $s\in(\mathcal V\cup\{\mathtt{M}\})^L$ contain committed tokens and
masks. The cached dLM produces logits $\ell_\theta(s)$. AdaBlock takes top-one proposals, commits
every masked position whose probability exceeds $\eta=0.9$, and always commits the most confident
remaining position. Denote this deterministic complete-state transition by
\begin{equation}
 F(s)=\operatorname{Transfer}_{\eta}\!\left(s,\ell_\theta(s)\right).
 \label{eq:transition}
\end{equation}
$F$ is a \emph{macro-transition}: it may reveal one or many tokens. The outer decoder chooses the
next block length from delimiter confidence and then repeatedly applies $F$ until the active block
contains no mask. Sequential AdaBlock therefore generates the trajectory
$s_0,s_1=F(s_0),\ldots,s_T=F^T(s_0)$.

\subsection{Draft, batch, verify}

After a verified transition, the already-computed logits provide top-one proposals and a confidence
ranking for the remaining masks. Starting at verified root $d_0=s_t$, we construct a cheap draft
chain $d_1,\ldots,d_D$ by revealing the next $w$ ranked proposals per level. No model call is used
for drafting. We then evaluate every node in one physical call,
\begin{equation}
 (\ell_0,\ldots,\ell_D)=
 f_\theta\!\left([d_0;\ldots;d_D],\operatorname{Repeat}(\mathrm{KV},D+1)\right).
 \label{eq:batch-verify}
\end{equation}
Verification begins at $j=0$. It computes the \emph{actual} AdaBlock successor
$u=\operatorname{Transfer}_\eta(d_j,\ell_j)$. If $u=d_{j+1}$ exactly as a complete token tensor,
the edge is accepted and verification continues at $j+1$. On the first mismatch, \method returns
$u$, not $d_{j+1}$. If all draft edges match, the verified leaf logits supply one additional native
transition. Thus one batched call realizes between one and $D+1$ sequential AdaBlock transitions.

\begin{proposition}[Trajectory equivalence]
Assume temperature zero and batch-invariant model outputs: evaluating a state alone or as a row of
Equation~\ref{eq:batch-verify} induces the same $F$. If a \method round starts at $s_t$, it returns
$s_{t+m}=F^m(s_t)$ for some $m\ge1$. Consequently, \method and sequential AdaBlock have identical
states, block boundaries, and generated tokens.
\end{proposition}

\noindent\textit{Proof sketch.}
The verified root equals $s_t$. An accepted edge is admitted only when its complete child equals
$F$ of the current verified state, so induction gives $d_j=F^j(s_t)$ along the accepted prefix. A
rejection returns the directly evaluated successor of the last verified node; a fully accepted chain
uses the directly evaluated leaf successor. In both cases the returned state is an iterate of $F$.
Applying the argument round-by-round and block-by-block proves the claim. Importantly, no assumption
about the draft-quality predictor appears in the proof.

\subsection{Risk controls capacity, never acceptance}

We reuse the compact local trajectory features that were predictive in failed early-stopping runs:
mask fraction, entropy and confidence quantiles, top-one margin, token churn, digit/delimiter rates,
and normalized temporal Jensen--Shannon change. A histogram-gradient-boosting model predicts whether
the current top-one block proposal differs from the native block endpoint. For risk score
$p_t=\widehat q(s_t)$, the router chooses
\begin{equation}
 D(p_t)=
 \begin{cases}
 D_{\max}, & p_t\le q_{\rm deep},\\
 D_{\rm med}, & q_{\rm deep}<p_t\le q_{\rm med},\\
 0, & p_t>q_{\rm med},
 \end{cases}
 \qquad
 w_t=\max\{1,\operatorname{round}(\gamma a_{t-1})\},
 \label{eq:router}
\end{equation}
where $a_{t-1}$ is the number of tokens revealed by the last verified transition. We predeclare
three routers $(D_{\max},D_{\rm med},q_{\rm deep},q_{\rm med},\gamma)$:
$(2,1,.05,.15,.75)$, $(4,2,.10,.30,1)$, and $(6,3,.20,.50,1)$. Prediction error only changes
batch size and acceptance opportunity. Exact token acceptance is always determined by
Equation~\ref{eq:transition}.

Physical NFE counts batched model invocations. We separately record the number of evaluated tree
nodes, accepted edges, verified native transitions, rejection depth, memory, and latency. This
prevents an NFE reduction obtained by a larger batch from being mistaken for a free reduction in
total FLOPs.

\section{Experimental protocol}

\paragraph{Models, data, and splits.}
We pin LLaDA-8B-Instruct and Dream-v0-Instruct-7B revisions and use deterministic decoding
with generation length 256, at most 64 refinement steps, and block length at most 64.
Estimator training, tuning, and calibration use disjoint deterministic subsets of GSM8K,
MATH, and MBPP training data. The 500-prompt calibration mixture per model is never used for
fitting or threshold search. The completed v1 evaluation is treated only as development
evidence. Workshop confirmation draws its deterministic complement: 500 GSM8K and 200 MATH-500,
100 sanitized MBPP, and all 64 remaining HumanEval problems; HumanEval is declared
out-of-domain. Proportional strata spread each sample across the unused benchmark rows. The
frozen YAML, dataset revisions, seeds, prior-sample exclusions, and protocol hash accompany
the code.

\paragraph{Comparators and ablations.}
The confirmatory comparison includes AdaBlock, the better fixed-depth verified speculator
($D=2$ or $D=4$) chosen on tuning data, and the frozen \method router. This directly tests whether
risk adaptation improves over the closest SSD-style nonlearned alternative while holding the
AdaBlock macro-verifier fixed. The local-risk router is evaluated with Brier score/AUROC on a
prompt-grouped holdout, while draft acceptance and rejection depth measure whether its compute
allocation is useful.

\paragraph{Metrics and inference.}
Primary metrics are mean NFE, exact task accuracy, elapsed time, harmful regressions, and
sequence disagreements.
Every method receives the same prompt order and decode settings. We report paired bootstrap
95\% intervals for NFE and accuracy differences (10,000 resamples), normalized cross-model
NFE, simultaneous calibration bounds, and a failure taxonomy. The positive headline is
allowed only if: (1) every paired generated token sequence equals AdaBlock and every record has
verified-transition evidence; (2) \method lowers physical NFE versus AdaBlock on both models and
versus the best nonlearned rule overall; and (3) the pooled in-domain paired-bootstrap NFE-reduction
lower bound exceeds 5\% for each model. We additionally report latency intervals and verifier-node
work; NFE alone is not presented as wall-clock speedup.

\paragraph{Compute-aware funnel.}
We first run 32 prompts/model and require exact generated-token, NFE-history, and block-history
parity between native AdaBlock and the instrumented loop. We reuse 600 native full-budget
traces/model from v4--v7 when available, but refit the v8 router and discard every earlier policy,
selection, certificate, and confirmation result. Three routers are screened on a fresh 150
prompts/model. Calibration on 500 prompts/model begins only when the selected router saves at
least 5\% physical NFE and beats the best eligible nonlearned control on \emph{each} model.
Calibration hard-fails on one sequence disagreement or one missing verifier audit. Confirmation
cannot start without this paired audit and an explicit flag. All records are atomic and resumable;
models are loaded one at a time on one A100.

\section{Results}

\IfFileExists{generated/headline.tex}{
  \noindentRC-PAG is not eligible for the positive headline under the predeclared gates.


  \begin{table}[t]
    \caption{Cross-model NFE summary on preregistered confirmation sets. Values are generated
    from paired, complete, non-mock records. Lower normalized NFE is better.}
    \label{tab:results}
    \centering
    \small
    \begin{tabular}{lrr}
\toprule
Variant & Normalized NFE & $\Delta$ vs. AdaBlock \\
\midrule
adablock & 1.000 & +0.0\% \\ 
best\_nonlearned & 0.968 & -3.2\% \\ 
rc\_pag\_history & 0.904 & -9.6\% \\ 
\bottomrule
\end{tabular}

  \end{table}

  \begin{figure}[t]
    \centering
    \includegraphics[width=0.72\linewidth]{generated/nfe_accuracy.pdf}
    \caption{Accuracy--compute frontier. Each point aggregates paired confirmatory examples;
    model/dataset detail is supplied in Appendix~\ref{app:results}.}
    \label{fig:frontier}
  \end{figure}
}{
  \pending{\textbf{Confirmatory results pending.} The build contains no generated numerical
  claims because a complete, non-mock confirmation directory was not supplied. The release
  script refuses mock artifacts. Table~\ref{tab:protocol} records the analysis that will be
  inserted without changing the method or claim gates.}
}

\begin{table}[t]
\caption{Frozen confirmatory decisions. These gates are evaluated mechanically before any
positive headline is emitted.}
\label{tab:protocol}
\centering
\small
\begin{tabular}{ll}
\toprule
Decision & Frozen rule \\
\midrule
Output safety & zero token-sequence disagreements with AdaBlock \\
Verification & every transition accepted only by exact state equality \\
Compute & model-wise paired-bootstrap reduction LCB $>5\%$ \\
Adaptive method & local risk chooses depth; native transition chooses tokens \\
Fallback & AdaBlock if exactness or compute gate fails \\
\bottomrule
\end{tabular}
\end{table}

Regardless of outcome, the final report includes all candidates and all datasets. Any exactness or
compute-gate failure triggers the conservative fallback; an NFE win without a latency win is
reported as reduced sequential depth, not as wall-clock acceleration.

\section{Limitations and broader impact}

The equivalence result assumes deterministic, batch-invariant transitions. Different batching
kernels can rarely perturb tied or nearly tied logits; we therefore audit exact tokens on every
calibration and confirmation prompt and treat any mismatch as a method failure. Losslessness is
relative to AdaBlock, not to a ground-truth answer, and provides no guarantee of correctness,
factuality, fairness, or harmlessness. Batched verification reduces sequential calls but evaluates
more states per call, so a physical-NFE win need not imply a latency or energy win; both are reported.
The current linear draft chain may underuse capacity when the number of threshold-qualified tokens
changes abruptly. Finally, two 7--8B model families and four benchmarks cannot establish universal
speedup. Faster generation can lower energy per response but may also increase aggregate use.

\section{Conclusion}

\method resolves the accuracy--compute conflict by removing acceptance authority from the learned
router. Risk allocates speculative capacity; the exact native AdaBlock transition remains the sole
token authority. This makes aggressive proposals compatible with trajectory equivalence and turns
prediction errors into efficiency errors. The frozen evaluation separately tests whether the
resulting reduction in sequential calls is large and fast enough to matter.

\typeout{RC-PAG-MAIN-PAGES:\arabic{page}}

\bibliographystyle{plainnat}
\bibliography{rc_pag_references}

\appendix

\section{Proof details}
\label{app:proof}

Let $r_0=s_t$. At verifier node $d_j=r_j$, the batched output induces exactly
$u_j=F(r_j)$. If $u_j=d_{j+1}$, the verifier sets $r_{j+1}=d_{j+1}=F(r_j)$; hence by induction
$r_j=F^j(s_t)$ along every accepted prefix. At the first rejection it returns
$u_j=F^{j+1}(s_t)$. If no edge rejects, applying the verified leaf output returns the same form.
Thus each round advances a positive integer number of reference transitions. Induction over rounds
finishes the active block at the same state. Since delimiter-based block sizing is evaluated only
at identical block boundaries, induction over blocks gives the complete generated-sequence result.

\section{Online algorithm}

For each block, \method executes the ordinary first AdaBlock transition and retains its logits and
immutable full-sequence cache. While masks remain it: (1) extracts local features and chooses depth
with Equation~\ref{eq:router}; (2) constructs the linear draft chain from stale top-one proposals;
(3) repeats the immutable cache and evaluates all nodes as one batch; (4) traverses only edges whose
complete child state equals the verified native successor; (5) returns the verified successor at a
mismatch or the verified leaf successor; and (6) records physical NFE, nodes, accepted edges,
rejection depth, and a sequence-safe flag. A depth-zero decision is exactly one native transition.

\section{Full confirmatory results}
\label{app:results}

\IfFileExists{generated/main_results.tex}{
  \begin{center}\scriptsize\begin{tabular}{lllrr}
\toprule
Model & Dataset & Method & Accuracy & Mean NFE \\
\midrule
llada & gsm8k\_test & adablock & 78.2 & 86.89 \\
llada & gsm8k\_test & best\_nonlearned & 78.2 & 86.89 \\
llada & gsm8k\_test & rc\_pag\_history & 78.2 & 85.06 \\
llada & math500 & adablock & 37.7 & 119.72 \\
llada & math500 & best\_nonlearned & 37.7 & 119.72 \\
llada & math500 & rc\_pag\_history & 37.3 & 118.17 \\
llada & mbpp\_sanitized & adablock & 4.0 & 71.66 \\
llada & mbpp\_sanitized & best\_nonlearned & 4.0 & 71.66 \\
llada & mbpp\_sanitized & rc\_pag\_history & 4.0 & 70.57 \\
llada & humaneval & adablock & 32.0 & 95.48 \\
llada & humaneval & best\_nonlearned & 32.0 & 95.48 \\
llada & humaneval & rc\_pag\_history & 33.0 & 94.39 \\
dream & gsm8k\_test & adablock & 69.8 & 103.04 \\
dream & gsm8k\_test & best\_nonlearned & 69.4 & 91.63 \\
dream & gsm8k\_test & rc\_pag\_history & 53.6 & 89.46 \\
dream & math500 & adablock & 45.3 & 123.56 \\
dream & math500 & best\_nonlearned & 41.0 & 121.44 \\
dream & math500 & rc\_pag\_history & 40.3 & 114.39 \\
dream & mbpp\_sanitized & adablock & 3.0 & 80.02 \\
dream & mbpp\_sanitized & best\_nonlearned & 4.0 & 74.69 \\
dream & mbpp\_sanitized & rc\_pag\_history & 0.0 & 54.25 \\
dream & humaneval & adablock & 48.0 & 84.78 \\
dream & humaneval & best\_nonlearned & 16.0 & 68.07 \\
dream & humaneval & rc\_pag\_history & 16.0 & 64.50 \\
\bottomrule
\end{tabular}
\end{center}
  \begin{center}\scriptsize\begin{tabular}{lrrrrr}
\toprule
Candidate & Errors/$n$ & Risk & Upper & $p$ & Certified \\
\midrule
dream/history\_t010\_p2 & 0/300 & 0.000 & 0.018 & 2.075e-07 & yes \\
dream/history\_t025\_p2 & 0/300 & 0.000 & 0.018 & 2.075e-07 & yes \\
dream/history\_t050\_p2 & 2/300 & 0.007 & 0.031 & 2.927e-05 & yes \\
dream/local\_t010\_p2 & 0/300 & 0.000 & 0.018 & 2.075e-07 & yes \\
dream/local\_t025\_p2 & 0/300 & 0.000 & 0.018 & 2.075e-07 & yes \\
dream/local\_t050\_p2 & 1/300 & 0.003 & 0.025 & 3.484e-06 & yes \\
llada/history\_t010\_p2 & 2/300 & 0.007 & 0.031 & 2.927e-05 & yes \\
llada/history\_t025\_p2 & 4/300 & 0.013 & 0.042 & 0.0006908 & yes \\
llada/history\_t050\_p2 & 5/300 & 0.017 & 0.047 & 0.002332 & yes \\
llada/local\_t010\_p2 & 2/300 & 0.007 & 0.031 & 2.927e-05 & yes \\
llada/local\_t025\_p2 & 4/300 & 0.013 & 0.042 & 0.0006908 & yes \\
llada/local\_t050\_p2 & 7/300 & 0.023 & 0.057 & 0.01597 & no \\
\bottomrule
\end{tabular}
\end{center}
  \begin{center}\scriptsize\begin{tabular}{llllrr}
\toprule
Model & Features & Estimator & Split & Brier & AUROC \\
\midrule
dream & rc\_pag\_history & hist\_gradient\_boosting & holdout & 0.026 & 0.990 \\
dream & rc\_pag\_history & logistic & holdout & 0.029 & 0.987 \\
dream & rc\_pag\_local & hist\_gradient\_boosting & holdout & 0.027 & 0.990 \\
dream & rc\_pag\_local & logistic & holdout & 0.029 & 0.987 \\
llada & rc\_pag\_history & hist\_gradient\_boosting & holdout & 0.051 & 0.982 \\
llada & rc\_pag\_history & logistic & holdout & 0.056 & 0.979 \\
llada & rc\_pag\_local & hist\_gradient\_boosting & holdout & 0.051 & 0.982 \\
llada & rc\_pag\_local & logistic & holdout & 0.056 & 0.979 \\
\bottomrule
\end{tabular}
\end{center}
  \IfFileExists{generated/screening_ablation.tex}{
    \begin{center}\scriptsize\input{generated/screening_ablation.tex}\end{center}
  }{}
  \begin{figure}[h]
    \centering
    \includegraphics[width=0.75\linewidth]{generated/reliability.pdf}
    \caption{Simultaneous calibration diagnostics.}
  \end{figure}
}{
  Complete tables are generated only from a validated non-mock run.
}

\section{Reproducibility checklist}

The artifact records repository and dataset revisions, the config and protocol hashes,
prompt IDs, raw outputs, correctness, NFE, elapsed time, per-block compact traces, selected
policy, parity and readiness audits, harmful regressions, raw paired NFE changes, harm calibration
$p$-values and simultaneous bounds, paired confirmation intervals, and every failed
claim gate. Code execution grading uses a separate subprocess with a timeout. Secrets are
read from the environment and never written to the command line or artifacts. Exact local
and Slurm commands are provided in \texttt{docs/rc\_pag\_runbook.md}.

\end{document}
